% GENERATED FILE. Edit canonical lessons, then run npm run generate:books.
% canonical-source-hash: fnv1a64:c481e823ce059603
% canonical-lessons: TE-C533-kapuram, TE-C533-samsaram, TE-C533-grhini, TE-C533-atithyam, TE-C533-muccatincu, TE-R533-review-526-533

\chapter{Home Life}
\label{ch:te-home-life}
\hlchaptermodality{533}

\begin{chapteropening}
\textbf{By the end of this chapter:} \emph{I can talk about married life, family life, hospitality and a good chat.}

Complete the last lesson of chapter 533: I can talk about married life, family life, hospitality and a good chat.
\end{chapteropening}

\section[\texorpdfstring{\te{కాపురం} (kāpuraṁ)}{కాపురం (kāpuraṁ)}]{\te{కాపురం} (kāpuraṁ) — married life; a household}

\label{lesson:TE-C533-kapuram}



\begin{quote}
Before the new one: say the Telugu for to coax, then the Telugu for a tantrum.
\end{quote}

\subsection*{You'll want to know: \te{కాపురం}}
\textbf{\te{కాపురం}} — \emph{kāpuraṁ} — \textquotedblleft{}married life; a household\textquotedblright{}.

\begin{quote}
\te{వాళ్ళ} \te{కాపురం} \te{సంతోషంగా} \te{సాగుతోంది}.
\end{quote}

\begin{quote}
\emph{vāḷḷa kāpuraṁ santōṣaṅgā sāgutōndi.}
\end{quote}

\begin{quote}
Their married life is going happily.
\end{quote}

\te{కాపురం} \te{పెట్టు}, to set up house.

\begin{cousinweb}
\textbf{\te{కాపురం}} (\emph{kāpuraṁ}) is native Dravidian: DEDR 1456 lists Telugu \emph{kāpuramu}, \textquotedblleft{}dwelling, residence\textquotedblright{}, beside \emph{kāpu}, \textquotedblleft{}a cultivator\textquotedblright{}.
\end{cousinweb}

\begin{grammarlens}[title={What you've built}]
One more word for describing what you lived and felt.
\end{grammarlens}

\subsection*{Your turn}
\begin{itemize}
\raggedright
  \item \emph{Say it:} \emph{kāpuraṁ}
  \item \emph{Say it:} \emph{kāpuraṁ}, once more
  \item \emph{Say it:} \emph{vāḷḷa kāpuraṁ santōṣaṅgā sāgutōndi.} — Their married life is going happily.
  \item \emph{Recall:} say the Telugu for to coax, then the Telugu for a tantrum, then say \emph{kāpuraṁ} again
\end{itemize}

\begin{tcolorbox}[breakable,colback=teal!4,colframe=teal!35!black,title={Before you move on}]
What does \te{కాపురం} mean? (\textquotedblleft{}Married life; a household\textquotedblright{}.) Say it once more.
\end{tcolorbox}
\section[\texorpdfstring{\te{సంసారం} (saṁsāraṁ)}{సంసారం (saṁsāraṁ)}]{\te{సంసారం} (saṁsāraṁ) — family life}

\label{lesson:TE-C533-samsaram}



\begin{quote}
Before the new one: say the Telugu for a tantrum, then the Telugu for married life.
\end{quote}

\subsection*{You'll want to know: \te{సంసారం}}
\textbf{\te{సంసారం}} — \emph{saṁsāraṁ} — \textquotedblleft{}family life\textquotedblright{}.

\begin{quote}
\te{పెద్ద} \te{సంసారం} \te{నడపడం} \te{సులభం} \te{కాదు}.
\end{quote}

\begin{quote}
\emph{pedda saṁsāraṁ naḍapaḍaṁ sulabhaṁ kādu.}
\end{quote}

\begin{quote}
Running a big family is not easy.
\end{quote}

Also one's family itself: \te{మా} \te{సంసారం}.

\begin{cousinweb}
\textbf{\te{సంసారం}} (\emph{saṁsāraṁ}) is borrowed from Sanskrit \emph{saṁsāra}, \textquotedblleft{}secular life, the world\textquotedblright{} (CDIAL 13036).
\end{cousinweb}

\begin{grammarlens}[title={What you've built}]
One more word for describing what you lived and felt.
\end{grammarlens}

\subsection*{Your turn}
\begin{itemize}
\raggedright
  \item \emph{Say it:} \emph{saṁsāraṁ}
  \item \emph{Say it:} \emph{saṁsāraṁ}, once more
  \item \emph{Say it:} \emph{pedda saṁsāraṁ naḍapaḍaṁ sulabhaṁ kādu.} — Running a big family is not easy.
  \item \emph{Recall:} say the Telugu for a tantrum, then the Telugu for married life, then say \emph{saṁsāraṁ} again
\end{itemize}

\begin{tcolorbox}[breakable,colback=teal!4,colframe=teal!35!black,title={Before you move on}]
What does \te{సంసారం} mean? (\textquotedblleft{}Family life\textquotedblright{}.) Say it once more.
\end{tcolorbox}
\section[\texorpdfstring{\te{గృహిణి} (gṛhiṇi)}{గృహిణి (gṛhiṇi)}]{\te{గృహిణి} (gṛhiṇi) — a homemaker}

\label{lesson:TE-C533-grhini}



\begin{quote}
Before the new one: say the Telugu for married life, then the Telugu for family life.
\end{quote}

\subsection*{You'll want to know: \te{గృహిణి}}
\textbf{\te{గృహిణి}} — \emph{gṛhiṇi} — \textquotedblleft{}a homemaker\textquotedblright{}.

\begin{quote}
\te{మా} \te{అమ్మ} \te{గృహిణి}.
\end{quote}

\begin{quote}
\emph{mā amma gṛhiṇi.}
\end{quote}

\begin{quote}
My mother is a homemaker.
\end{quote}

From \te{గృహం}, house.

\begin{cousinweb}
\textbf{\te{గృహిణి}} (\emph{gṛhiṇi}) is borrowed from Sanskrit \emph{gṛhiṇī}, \textquotedblleft{}the mistress of a house\textquotedblright{} (Monier-Williams).
\end{cousinweb}

\begin{grammarlens}[title={What you've built}]
One more word for describing what you lived and felt.
\end{grammarlens}

\subsection*{Your turn}
\begin{itemize}
\raggedright
  \item \emph{Say it:} \emph{gṛhiṇi}
  \item \emph{Say it:} \emph{gṛhiṇi}, once more
  \item \emph{Say it:} \emph{mā amma gṛhiṇi.} — My mother is a homemaker.
  \item \emph{Recall:} say the Telugu for married life, then the Telugu for family life, then say \emph{gṛhiṇi} again
\end{itemize}

\begin{tcolorbox}[breakable,colback=teal!4,colframe=teal!35!black,title={Before you move on}]
What does \te{గృహిణి} mean? (\textquotedblleft{}A homemaker\textquotedblright{}.) Say it once more.
\end{tcolorbox}
\section[\texorpdfstring{\te{ఆతిథ్యం} (ātithyaṁ)}{ఆతిథ్యం (ātithyaṁ)}]{\te{ఆతిథ్యం} (ātithyaṁ) — hospitality}

\label{lesson:TE-C533-atithyam}



\begin{quote}
Before the new one: say the Telugu for family life, then the Telugu for a homemaker.
\end{quote}

\subsection*{You'll want to know: \te{ఆతిథ్యం}}
\textbf{\te{ఆతిథ్యం}} — \emph{ātithyaṁ} — \textquotedblleft{}hospitality\textquotedblright{}.

\begin{quote}
\te{వాళ్ళ} \te{ఆతిథ్యం} \te{మాకు} \te{ఎంతో} \te{నచ్చింది}.
\end{quote}

\begin{quote}
\emph{vāḷḷa ātithyaṁ māku entō naccindi.}
\end{quote}

\begin{quote}
We liked their hospitality very much.
\end{quote}

Linked to \te{అతిథి}, a guest.

\begin{cousinweb}
\textbf{\te{ఆతిథ్యం}} (\emph{ātithyaṁ}) is borrowed from Sanskrit \emph{ātithya}, \textquotedblleft{}hospitable reception, hospitality\textquotedblright{} (Monier-Williams), built on \emph{atithi}, \textquotedblleft{}a guest\textquotedblright{}.
\end{cousinweb}

\begin{grammarlens}[title={What you've built}]
One more word for describing what you lived and felt.
\end{grammarlens}

\subsection*{Your turn}
\begin{itemize}
\raggedright
  \item \emph{Say it:} \emph{ātithyaṁ}
  \item \emph{Say it:} \emph{ātithyaṁ}, once more
  \item \emph{Say it:} \emph{vāḷḷa ātithyaṁ māku entō naccindi.} — We liked their hospitality very much.
  \item \emph{Recall:} say the Telugu for family life, then the Telugu for a homemaker, then say \emph{ātithyaṁ} again
\end{itemize}

\begin{tcolorbox}[breakable,colback=teal!4,colframe=teal!35!black,title={Before you move on}]
What does \te{ఆతిథ్యం} mean? (\textquotedblleft{}Hospitality\textquotedblright{}.) Say it once more.
\end{tcolorbox}
\section[\texorpdfstring{\te{ముచ్చటించు} (muccaṭiñcu)}{ముచ్చటించు (muccaṭiñcu)}]{\te{ముచ్చటించు} (muccaṭiñcu) — to chat; to have a heart-to-heart}

\label{lesson:TE-C533-muccatincu}



\begin{quote}
Before the new one: say the Telugu for a homemaker, then the Telugu for hospitality.
\end{quote}

\subsection*{You'll want to know: \te{ముచ్చటించు}}
\textbf{\te{ముచ్చటించు}} — \emph{muccaṭiñcu} — \textquotedblleft{}to chat; to have a heart-to-heart\textquotedblright{}.

\begin{quote}
\te{సాయంత్రం} \te{అమ్మమ్మ} \te{పక్కింటివాళ్ళతో} \te{ముచ్చటించింది}.
\end{quote}

\begin{quote}
\emph{sāyantraṁ ammamma pakkiṇṭivāḷḷatō muccaṭiñcindi.}
\end{quote}

\begin{quote}
In the evening grandma chatted with the neighbours.
\end{quote}

Past (I …): \te{ముచ్చటించాను}, \emph{muccaṭiñcānu}.

-\te{తో} \te{ముచ్చటించు}; \te{ముచ్చట}, a pleasant chat or a charming thing.

\begin{cousinweb}
\textbf{\te{ముచ్చటించు}} (\emph{muccaṭiñcu}) is built on \te{ముచ్చట}, \textquotedblleft{}a pleasant talk; a charming thing\textquotedblright{}, with the verb ending -\te{ఇంచు}; the origin of \te{ముచ్చట} is not settled.
\end{cousinweb}

\begin{grammarlens}[title={What you've built}]
One more word for describing what you lived and felt.
\end{grammarlens}

\subsection*{Your turn}
\begin{itemize}
\raggedright
  \item \emph{Say it:} \emph{muccaṭiñcu}
  \item \emph{Say it:} \emph{muccaṭiñcu}, once more
  \item \emph{Say it:} \emph{sāyantraṁ ammamma pakkiṇṭivāḷḷatō muccaṭiñcindi.} — In the evening grandma chatted with the neighbours.
  \item \emph{Recall:} say the Telugu for a homemaker, then the Telugu for hospitality, then say \emph{muccaṭiñcu} again
\end{itemize}

\begin{tcolorbox}[breakable,colback=teal!4,colframe=teal!35!black,title={Before you move on}]
What does \te{ముచ్చటించు} mean? (\textquotedblleft{}To chat; to have a heart-to-heart\textquotedblright{}.) Say it once more.
\end{tcolorbox}
\section[(dialogue)]{Words from chapters 526–533 — a review}

\label{lesson:TE-R533-review-526-533}



\begin{quote}
Say the Telugu for to appoint and to manage.
\end{quote}

\subsection*{Your turn}
Say these aloud:

\begin{itemize}
\raggedright
  \item to appoint, to manage, to supervise, to cooperate, training
  \item to achieve, to excel, to fail, a resignation, to quit
  \item a sibling, brothers, a husband's younger brother, a husband's sister, an elder brother's wife
  \item a maternal uncle, a paternal aunt, a father's elder brother, a mother's elder sister, a niece
  \item a co-father-in-law, a great-grandfather, a nephew, a kinsman, a granny
  \item an infant, a newborn, offspring, an heir, a lullaby
  \item upbringing, pampering, to support, to coax, a tantrum
  \item married life, family life, a homemaker, hospitality, to chat
\end{itemize}

\begin{tcolorbox}[breakable,colback=teal!4,colframe=teal!35!black,title={Before you move on}]
To appoint? (\textbf{\te{నియమించు}}, \emph{niyamiñcu}.) To manage? (\textbf{\te{నిర్వహించు}}, \emph{nirvahiñcu}.) To supervise? (\textbf{\te{పర్యవేక్షించు}}, \emph{paryavēkṣiñcu}.) To cooperate? (\textbf{\te{సహకరించు}}, \emph{sahakariñcu}.) Training? (\textbf{\te{శిక్షణ}}, \emph{śikṣaṇa}.) To achieve? (\textbf{\te{సాధించు}}, \emph{sādhiñcu}.) To excel? (\textbf{\te{రాణించు}}, \emph{rāṇiñcu}.) To fail? (\textbf{\te{విఫలమవు}}, \emph{viphalamavu}.) A resignation? (\textbf{\te{రాజీనామా}}, \emph{rājīnāmā}.) To quit? (\textbf{\te{మానేయు}}, \emph{mānēyu}.) A sibling? (\textbf{\te{తోబుట్టువు}}, \emph{tōbuṭṭuvu}.) Brothers? (\textbf{\te{అన్నదమ్ములు}}, \emph{annadammulu}.) A husband's younger brother? (\textbf{\te{మరిది}}, \emph{maridi}.) A husband's sister? (\textbf{\te{ఆడపడుచు}}, \emph{āḍapaḍucu}.) An elder brother's wife? (\textbf{\te{వదిన}}, \emph{vadina}.) A maternal uncle? (\textbf{\te{మేనమామ}}, \emph{mēnamāma}.) A paternal aunt? (\textbf{\te{మేనత్త}}, \emph{mēnatta}.) A father's elder brother? (\textbf{\te{పెద్దనాన్న}}, \emph{peddanānna}.) A mother's elder sister? (\textbf{\te{పెద్దమ్మ}}, \emph{peddamma}.) A niece? (\textbf{\te{మేనకోడలు}}, \emph{mēnakōḍalu}.) A co-father-in-law? (\textbf{\te{వియ్యంకుడు}}, \emph{viyyaṅkuḍu}.) A great-grandfather? (\textbf{\te{ముత్తాత}}, \emph{muttāta}.) A nephew? (\textbf{\te{మేనల్లుడు}}, \emph{mēnalluḍu}.) A kinsman? (\textbf{\te{చుట్టం}}, \emph{cuṭṭaṁ}.) A granny? (\textbf{\te{అవ్వ}}, \emph{avva}.) An infant? (\textbf{\te{బిడ్డ}}, \emph{biḍḍa}.) A newborn? (\textbf{\te{శిశువు}}, \emph{śiśuvu}.) Offspring? (\textbf{\te{సంతానం}}, \emph{santānaṁ}.) An heir? (\textbf{\te{వారసుడు}}, \emph{vārasuḍu}.) A lullaby? (\textbf{\te{జోలపాట}}, \emph{jōlapāṭa}.) Upbringing? (\textbf{\te{పెంపకం}}, \emph{pempakaṁ}.) Pampering? (\textbf{\te{గారాబం}}, \emph{gārābaṁ}.) To support? (\textbf{\te{పోషించు}}, \emph{pōṣiñcu}.) To coax? (\textbf{\te{బుజ్జగించు}}, \emph{bujjagiñcu}.) A tantrum? (\textbf{\te{మారాం}}, \emph{mārāṁ}.) Married life? (\textbf{\te{కాపురం}}, \emph{kāpuraṁ}.) Family life? (\textbf{\te{సంసారం}}, \emph{saṁsāraṁ}.) A homemaker? (\textbf{\te{గృహిణి}}, \emph{gṛhiṇi}.) Hospitality? (\textbf{\te{ఆతిథ్యం}}, \emph{ātithyaṁ}.) To chat? (\textbf{\te{ముచ్చటించు}}, \emph{muccaṭiñcu}.)
\end{tcolorbox}
